../cictikz/src/cictikz/data/tex/cictikz_preamble.tex

% Everything between a bond pad and the first gate of a digital input.
%
% Read it left to right, because that is the order the current sees it.
% The pad's two diodes take a zap to whichever rail it belongs to. The
% supply clamp - a grounded gate NMOS and a diode, the same pair as the
% network figure earlier in the chapter - is what makes those diodes
% useful, since a diode to VDD only helps if VDD itself has somewhere to
% dump the charge. The series resistor then limits what is left, and the
% secondary clamps hold the buffer's gate oxide at a rail while the
% primary devices, which are large and slow, get around to turning on.
%
% Two resistors and two stages, and that is the point of the figure: the
% primary devices are sized for amps and clamp late, the secondary ones
% are small and clamp early, and the resistor between them is what lets
% the two disagree.
%
% Both secondary devices are a transistor with its gate tied to the rail
% its channel terminal sits on, so both are off in normal operation.
%
% Differences from the hand drawing this replaces: the original marks
% the block boundaries with small open circles on the rails, which are
% ports on the sheet it was copied from and would read as junctions
% here; and it hops the signal line over the clamp column, where the
% house convention is a plain crossing.

\begin{circuitikz}[american, thick, transform shape, circuit ee IEC]

  ../cictikz/src/cictikz/data/tex/cictikz_lib.tex
  %%%%%%%%%%%%%%%%%%%%%%%%%%%%%%%%%%%%%%%%%%%%%%%%%%%%%%%%%%%%%%%%%%%%%%
%% Shared frame for the ESD build-up sequence in l02_esd.
%%
%% Four figures in the chapter are the same picture with more added each
%% time: a chip outline with three rails brought out to pads, then one
%% diode, then two, then the whole network with the grounded gate NMOS.
%% They only teach anything if they are literally the same drawing, so
%% the frame lives here and each figure adds its own devices.
%%
%% Coordinates: pads at x = 0, rails run to \esdRight. VDD at \esdVdd,
%% the signal pin at \esdPin, VSS at zero.
%%%%%%%%%%%%%%%%%%%%%%%%%%%%%%%%%%%%%%%%%%%%%%%%%%%%%%%%%%%%%%%%%%%%%%

\newcommand{\esdVdd}{5.0}
\newcommand{\esdPin}{2.6}
\newcommand{\esdRight}{8.8}

% A bond pad: the crossed square the chapter's originals use.
\newcommand{\esdPad}[1]{%
  \draw (-0.17,{#1-0.17}) rectangle (0.17,{#1+0.17});
  \draw (-0.17,{#1-0.17}) -- (0.17,{#1+0.17});
  \draw (-0.17,{#1+0.17}) -- (0.17,{#1-0.17});
}

% The chip outline, three rails, their pads and their names. #1 is the
% name of the middle rail, because the overview figure calls it PIN and
% the protection figures treat it as one signal pin among many.
\newcommand{\esdframe}{%
  \draw[dotted, gray] (-0.8,-1.3) rectangle ({\esdRight+0.3},{\esdVdd+1.3});
  \foreach \y/\name/\num in {\esdVdd/VDD/1, \esdPin/PIN/2, 0/VSS/0} {
    \draw (0,\y) -- (\esdRight,\y);
    \esdPad{\y}
    \node[armygreen, anchor=south west, scale=0.85] at (0.05,{\y+0.2}) {\name};
    %- above the rail, not left of it: the zap wiring comes in on the
    %- left and the numbers were landing on top of it
    \node[red, anchor=south east, scale=0.9] at (-0.25,{\y+0.18}) {\num};
  }
}

% A vertical diode conducting upward, from (#1,#2) to (#1,#3), with its
% body centred at #4 and named #5. Every protection device in these
% figures is one of these; drawing them through one macro is what makes
% the reverse-biased-in-normal-operation argument in the text hold
% across all four pictures.
%
% The centre is explicit because a diode spanning VSS to VDD would
% otherwise put its symbol exactly on the pin rail it is meant to cross
% without connecting.
\newcommand{\esdDiode}[5]{%
  \draw (#1,#2) -- (#1,{#4-0.45});
  \draw (#1,{#4-0.45}) to [D] (#1,{#4+0.45});
  \draw (#1,{#4+0.45}) -- (#1,#3);
  \fill (#1,#2) circle (0.075);
  \fill (#1,#3) circle (0.075);
  \node[red, anchor=west, scale=0.8] at ({#1+0.32},#4) {#5};
}

% A grounded gate NMOS at x=#1, drain on the rail at #2, source on the
% rail at #3, body centred at #4, node id #5 (letters only - it names a
% TikZ node) and caption #6. The centre is explicit for the same reason
% the diode's is: left to the midpoint, a VDD to VSS device lands on the
% pin rail it is meant to cross without touching. The gate goes to the source through a resistor,
% which is the whole trick: the resistor does nothing at DC, so the
% device is off whenever the chip is running, but during a zap the
% substrate current it develops lifts the gate and turns the parasitic
% bipolar on.
\newcommand{\esdggn}[6]{%
  \node[nmos] (esdm#5) at ({#1+0.35},#4) {};
  \draw (esdm#5.D) -- ({#1+0.35},#2);
  \draw (esdm#5.S) -- ({#1+0.35},#3);
  \fill ({#1+0.35},#2) circle (0.075);
  \fill ({#1+0.35},#3) circle (0.075);
  \draw (esdm#5.G) -- ({#1-0.35},#4);
  \draw ({#1-0.35},#4) to [R] ({#1-0.35},{#3+0.5});
  \draw ({#1-0.35},{#3+0.5}) -- ({#1-0.35},#3) -- ({#1+0.35},#3);
  %- beside the device, not above it: above puts it on the rail the
  %- device crosses without connecting
  \node[red, anchor=west, scale=0.8, align=left]
    at ({#1+0.95},#4) {#6};
}

% The zap itself: current forced in at one pad, taken out at another.
\newcommand{\esdZapIn}[1]{%
  \draw (-2.7,#1) to [I] (-1.3,#1);
  \draw (-1.3,#1) -- (0,#1);
}
\newcommand{\esdZapOut}[1]{%
  \draw (0,#1) -- (-1.3,#1);
  \draw (-1.3,#1) \vground;
}


  \def\yvdd{4.6}
  \def\ysig{2.6}
  \def\yclamp{1.7}
  \def\ygnd{0.2}

  %%- the pad, and what it is connected to the world with
  \draw (-0.45,{\ysig-0.45}) rectangle (0.45,{\ysig+0.45});
  \draw (-0.45,{\ysig-0.45}) -- (0.45,{\ysig+0.45});
  \draw (-0.45,{\ysig+0.45}) -- (0.45,{\ysig-0.45});
  \node[anchor=north] at (0,{\ysig-0.65}) {Pad};
  \draw[gray, dotted, very thick]
    (-2.6,{\yvdd+0.5}) .. controls (-1.6,{\yvdd+0.2}) and (-0.5,4.0) .. (0,{\ysig+0.3});
  \node[anchor=south] at (-1.9,{\yvdd+0.75}) {Bondwire};
  \draw (0.45,\ysig) -- (2.6,\ysig);

  %%- primary: a diode to each rail at the pad itself
  \draw (2.6,\yvdd) -- (8.0,\yvdd);
  \draw (2.6,\yvdd) \vsupply;
  \esdDiode{2.6}{\ysig}{\yvdd}{3.6}{}
  %- drawn out rather than through \esdDiode, which dots both of its
  %- ends: the bottom end here is a ground symbol, not a junction
  \draw (2.6,\ygnd) \vground;
  \draw (2.6,\ygnd) -- (2.6,0.6) \vdiode{} -- (2.6,\ysig);
  \fill (2.6,\ysig) circle (0.075);

  %%- and the resistor that separates the two stages
  \draw (2.6,\ysig) -- (3.0,\ysig);
  \draw (3.0,\ysig) \hresistor{};
  \draw (4.6,\ysig) -- (8.4,\ysig);

  %%- the supply clamp: the pad diodes are only as good as this is
  \draw (5.6,\yvdd) -- (5.6,\yclamp);
  \draw (5.0,\yclamp) -- (7.0,\yclamp);
  \fill (5.6,\yclamp) circle (0.075);
  \fill (5.6,\yvdd) circle (0.075);

  \node[nmos] (mclamp) at (5.0,1.0) {};
  \draw (mclamp.D) -- (5.0,\yclamp);
  \draw (mclamp.S) -- (5.0,\ygnd);
  \draw (5.0,\ygnd) \vground;
  \draw (mclamp.G) -- (4.2,1.0) -- (4.2,\ygnd);
  \draw (4.2,\ygnd) \vground;

  \draw (7.0,\ygnd) \vground;
  \draw (7.0,\ygnd) -- (7.0,0.4) \vdiode{} -- (7.0,\yclamp);
  \fill (7.0,\yclamp) circle (0.075);

  %%- secondary: a smaller resistor, and a clamp to each rail
  \draw (8.4,\ysig) \hresistor{};
  \draw (10.0,\ysig) -- (11.6,\ysig);
  \draw (8.0,\yvdd) -- (12.6,\yvdd);

  \node[nmos, yscale=-1] (mtop) at (11.0,3.6) {};
  \draw (mtop.S) -- (11.0,\yvdd);
  \fill (11.0,\yvdd) circle (0.075);
  \draw (mtop.D) -| (11.6,\ysig);
  \draw (mtop.G) -- (10.2,3.6) -- (10.2,\yvdd);
  \fill (10.2,\yvdd) circle (0.075);

  \node[nmos] (mbot) at (11.0,1.0) {};
  \draw (mbot.D) -| (11.6,\ysig);
  \draw (mbot.S) -- (11.0,\ygnd);
  \draw (11.0,\ygnd) \vground;
  \draw (mbot.G) -- (10.2,1.0) -- (10.2,\ygnd);
  \draw (10.2,\ygnd) \vground;

  \fill (11.6,\ysig) circle (0.075);

  %%- and what all of it is for
  \draw[rounded corners=3pt] (12.6,0.6) rectangle (14.6,\yvdd+0.6);
  \node[anchor=center, align=center] at (13.6,{\ysig+0.3}) {Input\\Buffer};
  \draw (11.6,\ysig) -- (12.6,\ysig);

  %%- the two stages, named
  \draw[armygreen, dotted, very thick] (2.0,-0.7) rectangle (8.2,5.9);
  \node[anchor=north] at (5.1,5.75) {Primary ESD};
  \draw[armygreen, dotted, very thick] (8.6,-0.3) rectangle (12.2,5.9);
  \node[anchor=north, align=center] at (10.4,5.75) {Secondary\\ESD};

\end{circuitikz}
\end{document}
